\hmsubsection{Development Methodology}{OAH}{}

The project was carried out using an iterative and incremental development approach 
based on the Scrum framework \cite{schwaber2020scrum}. This methodology was selected to enable continuous 
development, frequent feedback, and adaptability throughout the project lifecycle.

Work was organised into sprints with a duration of one to two weeks. Each sprint 
consisted of planning, implementation, and evaluation phases. At the end of each 
sprint, a retrospective was conducted to identify what worked well, address 
challenges, and define improvements for the next iteration. One notable adaptation 
during the project was the early decision to abandon MobileNetV2 transfer learning 
in favour of a classical computer vision pipeline, following poor real-time 
performance on the target hardware.

Task planning and tracking were managed using a digital Trello board, with tasks 
categorised into \textit{Backlog}, \textit{In Progress}, \textit{Testing}, and 
\textit{Done}. Stand-up meetings were held three times per week on Tuesdays, 
Wednesdays, and Thursdays, used to synchronise team members and identify blockers.